% Fonctionnement d'une base de données vectorielle
\begin{tikzpicture}[node distance=5mm and 9mm]

  % --- Indexation ---
  \node[btptitleband, text width=11.5cm] (b1) {Indexation};
  \node[btpbox, text width=2.3cm, below left=6mm and 3.2cm of b1.south] (doc) {Documents};
  \node[btpbox, text width=2.5cm, right=of doc] (emb1) {Modèle de\\plongement};
  \node[btpbox, text width=2.5cm, right=of emb1] (vec1) {Vecteurs\\$\mathbb{R}^{d}$};
  \node[btpdb, right=of vec1, minimum width=2.2cm] (store) {Index\\vectoriel};

  \draw[btparrow] (doc) -- (emb1);
  \draw[btparrow] (emb1) -- (vec1);
  \draw[btparrow] (vec1) -- (store);

  % --- Recherche ---
  \node[btptitleband, text width=11.5cm, below=12mm of doc.south west, anchor=north west] (b2) {Recherche};
  \node[btpboxalt, text width=2.3cm, below=6mm of b2.south west, anchor=north west] (query) {Question};
  \node[btpboxalt, text width=2.5cm, right=of query] (emb2) {Modèle de\\plongement};
  \node[btpboxalt, text width=2.9cm, right=of emb2] (knn) {Plus proches\\voisins\\[-1pt]\btplabel{HNSW, IVF, PQ}};
  \node[btpboxstrong, text width=2.4cm, right=of knn] (rank) {Classement par\\similarité};

  \draw[btparrow] (query) -- (emb2);
  \draw[btparrow] (emb2) -- (knn);
  \draw[btparrow] (knn) -- (rank);
  \draw[btparrowlight] (store.south) -- ++(0,-0.5) -| (knn.north);

  \node[btplabel, below=4mm of knn, text width=6.5cm]
    {Distance cosinus ou euclidienne, calculée de manière approchée
     pour rester sous-linéaire};

\end{tikzpicture}
